\documentclass[10pt]{article}
\usepackage[a4paper,margin=22mm]{geometry}
\usepackage{amsmath,amssymb,booktabs,xcolor}
\setlength{\parskip}{0.35em}\setlength{\parindent}{0pt}
\newcommand{\eqrefX}[1]{(\ref{#1})}
\makeatletter\def\eqref#1{(\ref{#1})}\makeatother
\newcommand{\pl}[1]{\textcolor{blue}{(#1)}}
\begin{document}
\providecommand{\labelX}{}
\def\eqpivotfree{}
% ---------------------------------------------------------------------
% Draft of Sec. VIII (Results) for the IEEE Access manuscript.
%
% Every number below is produced by a script in analysis/ and is marked
% with it in a comment, so each can be regenerated:
%   VIII-A  analysis/error_budget.py           (gate-passing counts)
%   VIII-B  analysis/onset_benchmark.py        (CV, detector comparison)
%   VIII-C  analysis/offset_accuracy_stats.py  (accuracy, per configuration)
%           analysis/offset_ge_variants.py     (ground-effect variants)
%   VIII-D  analysis/offset_accuracy_stats.py  (variance decomposition)
%           analysis/error_budget.py           (propagated budget)
%           analysis/vi_e_geometric_box.py     (a priori guarantee)
%   VIII-E  analysis/moff_flight_check.py      (flown vs final pipeline)
%
% PLACEHOLDERS the author must fill -- marked \TODO below:
%   - the excitation protocol counts (runs per configuration per direction)
%   - the free-flight tracking metrics with and without compensation
%     (this file has the commanded offsets, not the flight performance)
%
% Cross-references assumed to exist in the host document; rename before
% compiling into the manuscript:
%   \eqref{eq:cosh}      the two-segment rate model
%   \pl{weight eq}    weight from the moments
%   \pl{offset eq}    offset per direction
%   \pl{pivot-free avg} pivot-free average
%   Sec.~VI-E   the deviation bound
% ---------------------------------------------------------------------

\newcommand{\TODO}[1]{\textcolor{red}{[\textbf{TODO:} #1]}}

\section*{VIII. Results}

\subsection*{Excitation campaign and gating}

Five centre-of-mass configurations were produced by relocating ballast
within the planar offset box, giving ten independent case/axis
configurations once roll and pitch are counted separately. Each was
excited in both tip directions at commanded ramp rates of $0.10$,
$0.30$, $0.65$ and $1.20$~N\,m/s. \TODO{state the number of runs per
configuration per direction and the total attempted}

Runs enter the identification only through the gates of Sec.~VI-B: a
commanded-versus-realised rate error below $3\%$, a linearity residual
below $30$~mN\,m RMSE about the best-fit line, and a minimum
post-onset sample count. $140$ runs passed and are the basis of
everything below. The gates are deliberately blunt --- they reject
ramps that were stepped or aborted rather than tuning the fit --- and
no run was removed on the basis of its identified value.

\subsection*{Onset repeatability and the choice of detector}

Within a configuration and direction, the identified onset moment
repeats to a coefficient of variation of $2.9\%$ on average and
$5.8\%$ at worst. That figure is the primary measure of the method's
precision, and it is what the deviation bound of Sec.~VI-E
predicts should be achievable.

Because a two-segment fit could in principle be finding structure that
a simpler detector would find too, six alternatives were run over the
identical gated set and carried through the identical inversion
(Table~\ref{tab:bench}). Two are variants of the proposed estimator
--- one using CAD rather than identified constants, one widening the
search box --- and four are model-agnostic: a per-run nonlinear least
squares, two change-point detectors, and a CUSUM.

\begin{table}[t]
\centering
\caption{Onset detectors over the same $140$ gated runs, each carried
through the same inversion. ``CV'' is the within-configuration
coefficient of variation of the onset moment; ``offset RMS'' is the
delivered quantity against the load cell. $\Delta|$err$|$ is the mean
paired change in $|$error$|$ against the proposed method, with an
exact sign test over the ten configurations.}
\label{tab:bench}
\small
\setlength{\tabcolsep}{4pt}
\begin{tabular}{@{}l rr rr rr@{}}
\hline
 & CV & CV & offset & mean & $\Delta|$err$|$ & $p$ \\
method & mean & worst & RMS & & & \\
 & [\%] & [\%] & [mm] & [mm] & [mm] & \\
\hline
COSH (proposed)      & $2.9$ & $5.8$  & $1.65$ & $-0.93$ & ---     & --- \\
COSH, CAD constants  & $4.4$ & $12.3$ & $1.64$ & $-0.86$ & $+0.00$ & $0.75$\\
COSH, wide box       & $2.9$ & $5.7$  & $1.64$ & $-0.95$ & $-0.02$ & $0.75$\\
per-run NLS          & $6.7$ & $19.8$ & $1.72$ & $-1.07$ & $+0.02$ & $0.34$\\
CPD, Gaussian cost   & $4.2$ & $8.6$  & $1.67$ & $-0.91$ & $-0.03$ & $0.75$\\
CPD, RBF kernel      & $9.3$ & $27.4$ & $1.65$ & $-0.46$ & $-0.09$ & $1.00$\\
CUSUM, mean change   & $7.7$ & $24.0$ & $1.82$ & $-0.58$ & $+0.29$ & $0.75$\\
\hline
\end{tabular}
\end{table}

Two things follow, and they point in opposite directions. The proposed
estimator is the most repeatable by a wide margin --- $2.9\%$ against
$4.2$ to $9.3\%$ --- which matters for a pre-flight check that must
return the same answer twice. But every method delivers an offset RMS
between $1.64$ and $1.82$~mm, no paired comparison approaches
significance ($p\ge0.34$), and the delivered offsets agree to
$1.09$~mm in the median across all seven. The onset is therefore
\emph{not} the limiting step: the accuracy of the deliverable is set
elsewhere, which Sec.~VIII-D takes up.

\subsection*{Identified offset against the load cell}

The deliverable is the pivot-free average \pl{pivot-free avg},
inverted to a centre-of-mass offset and compared against an
independent load-cell measurement of the same quantity
(Table~\ref{tab:offset}).

\begin{table}[t]
\centering
\caption{Identified centre-of-mass offset against the load-cell truth,
over the ten configurations. ``repeat sd'' is the standard deviation
over run pairs within the configuration, and ``err/sd'' expresses the
error in units of that scatter.}
\label{tab:offset}
\small
\setlength{\tabcolsep}{4pt}
\begin{tabular}{@{}ll rrr rr@{}}
\hline
case & comp. & ident. & truth & error & repeat & err \\
 & & [mm] & [mm] & [mm] & sd [mm] & /sd \\
\hline
case\_01 & $y_{\rm off}$ & $-1.60$  & $-2.90$  & $+1.30$ & $1.09$ & $1.2$\\
case\_01 & $x_{\rm off}$ & $-11.18$ & $-11.45$ & $+0.27$ & $0.47$ & $0.6$\\
case\_02 & $y_{\rm off}$ & $-15.15$ & $-14.29$ & $-0.86$ & $0.40$ & $2.1$\\
case\_02 & $x_{\rm off}$ & $-10.57$ & $-9.90$  & $-0.67$ & $0.44$ & $1.5$\\
case\_03 & $y_{\rm off}$ & $-8.60$  & $-5.26$  & $-3.34$ & $0.51$ & $6.5$\\
case\_03 & $x_{\rm off}$ & $0.72$   & $3.14$   & $-2.42$ & $0.44$ & $5.4$\\
case\_04 & $y_{\rm off}$ & $5.24$   & $6.67$   & $-1.43$ & $0.74$ & $1.9$\\
case\_04 & $x_{\rm off}$ & $0.22$   & $2.40$   & $-2.18$ & $0.69$ & $3.2$\\
case\_05 & $y_{\rm off}$ & $10.79$  & $10.91$  & $-0.12$ & $0.82$ & $0.2$\\
case\_05 & $x_{\rm off}$ & $-10.57$ & $-10.89$ & $+0.32$ & $0.77$ & $0.4$\\
\hline
\end{tabular}
\end{table}

Over the ten configurations the RMS error is $1.64$~mm, with a
bootstrap $95\%$ confidence interval of $0.92$ to $2.24$~mm ($20$k
resamples), and the worst configuration is $3.34$~mm. Against offsets
that span $\pm15$~mm this is a relative accuracy of roughly one part
in ten of the range identified.

The mean error is $-0.91$~mm with a confidence interval of $-1.94$ to
$+0.11$~mm; seven of ten errors are negative, which an exact sign test
does not establish ($p = 0.34$). We therefore report the negative mean
as a point estimate and not as a demonstrated bias --- ten
configurations do not settle it, and the interval includes zero.

Two internal consistency checks accompany the comparison, neither of
which uses the load cell. Eq.~\pl{offset eq} estimates the same
offset twice, once from each tip direction with that direction's own
contact arm, and \pl{weight eq} returns the vehicle weight from the
onset moments. Without a ground-effect term the two directional
estimates disagree by $12.6$~mm in the median and \pl{weight eq}
returns a weight $5.2\%$ below $mg$ in all ten configurations;
restoring the interference channels brings the spread to $3.6$~mm and
the weight to $1.1\%$ below. The identified offset itself is reported
without a ground-effect model, to match the quantity that is flown;
Sec.~VI-B bounds what that omission costs at about half a millimetre
either way.

\subsection*{What limits the accuracy}

The errors of Table~\ref{tab:offset} are not run scatter. Within a
configuration the run-pair standard deviation is $0.60$~mm in the
median and $1.09$~mm at worst, whereas the standard deviation of the
error \emph{between} configurations is $1.43$~mm --- a ratio of
$2.4$. Read against each configuration's own scatter, three of the ten
errors exceed three standard deviations, the worst reaching $6.5$.
The error is therefore systematic per configuration, and repeating a
configuration cannot reduce it. This is the central statistical
statement of the paper: the reported accuracy is limited by the
ground-contact geometry, not by the number of excitation runs.

The estimator's own contribution is smaller by orders of magnitude.
Propagating the measured out-of-span forcing through the exact kernel
(Sec.~VI-E) gives a residual of $2.75$~mN\,m in the median
($12.28$ at worst) on the gravity channel and $0.53$ ($3.20$) on the
ground-effect channel, a rate deviation of $0.24\%$ of the nominal in
the median, and a threshold error of $0.04$~mN\,m in the median and
$0.15$ at worst --- $0.001$ to $0.004$~mm of offset. The a priori
guarantee over the whole admissible operating box, which assumes
nothing about the runs beyond the tilt cap and the ramp-rate range,
is $0.174$~mm for roll and $0.140$~mm for pitch. Either figure is an
order of magnitude below the $1.64$~mm actually achieved, which
locates the residual outside the estimator and is consistent with the
detector comparison of Table~\ref{tab:bench}.

\subsection*{Free-flight compensation}

The identified offsets were applied as a feed-forward moment trim at
take-off. The values flown are compared in Table~\ref{tab:flown}
against what the final pipeline returns, since the pipeline was
refined after the flights were performed.

\begin{table}[t]
\centering
\caption{Moment offsets commanded in the free-flight tests against the
final pipeline. ``std'' is the scatter of the identification at the
time of flight.}
\label{tab:flown}
\small
\setlength{\tabcolsep}{4pt}
\begin{tabular}{@{}ll rrr r@{}}
\hline
case & comp. & flown & std & final & diff \\
 & & [N\,m] & [N\,m] & [N\,m] & /std \\
\hline
case\_01 & $M_{{\rm off},x}$ & $-0.064$ & $0.037$ & $-0.048$ & $0.43$\\
case\_01 & $M_{{\rm off},y}$ & $0.331$  & $0.016$ & $0.336$  & $0.33$\\
case\_02 & $M_{{\rm off},x}$ & $-0.485$ & $0.009$ & $-0.479$ & $0.71$\\
case\_02 & $M_{{\rm off},y}$ & $0.333$  & $0.010$ & $0.334$  & $0.10$\\
case\_03 & $M_{{\rm off},x}$ & $-0.282$ & $0.028$ & $-0.272$ & $0.37$\\
case\_03 & $M_{{\rm off},y}$ & $-0.006$ & $0.019$ & $-0.023$ & $0.89$\\
case\_04 & $M_{{\rm off},x}$ & $0.173$  & $0.031$ & $0.165$  & $0.24$\\
case\_04 & $M_{{\rm off},y}$ & $-0.012$ & $0.024$ & $-0.007$ & $0.21$\\
case\_05 & $M_{{\rm off},x}$ & $0.307$  & $0.026$ & $0.341$  & $1.30$\\
case\_05 & $M_{{\rm off},y}$ & $0.343$  & $0.031$ & $0.334$  & $0.30$\\
\hline
\end{tabular}
\end{table}

The two agree to $8$~mN\,m in the median and $34$~mN\,m at worst, and
nine of the ten components fall inside the identification scatter that
was present at the time of flight; the worst is $1.3$ standard
deviations. Re-flying would therefore command the same moments to
within the precision of the identification itself, and the free-flight
results are reported against the pipeline as described.

\TODO{Report the flight performance here: attitude tracking error and
take-off transient with and without compensation, three repetitions per
configuration per condition, with the paired statistic across the five
configurations. This file supplies the commanded offsets only.}

\subsection*{Summary}

The method identifies a centre-of-mass offset to $1.64$~mm RMS
($95\%$ CI $0.92$--$2.24$) without a load cell, a test stand, or any
flight, from ground-contact excitation alone. It repeats to $2.9\%$ on
the onset moment, better than every model-agnostic baseline tried,
although the delivered offset is insensitive to that choice. Its own
modelling error is bounded a priori at $0.174$~mm and measured at
$0.001$--$0.004$~mm, so what limits the accuracy is the contact
geometry rather than the estimator or the number of repetitions ---
a conclusion the variance decomposition establishes directly, and one
that determines where further work should go.
\end{document}
